\documentclass[10pt]{extarticle}
\usepackage[a4paper,margin=1.5cm]{geometry}
\usepackage[T1]{fontenc}
\usepackage{amsmath,amssymb}
\usepackage{graphicx}
\usepackage[export]{adjustbox}
\usepackage{multicol}
\usepackage{enumitem}
\usepackage{booktabs,array,tabularx}
\setlength{\columnsep}{0.7cm}
\setlength{\parindent}{0pt}
\setlength{\parskip}{3pt}
\pagestyle{empty}
\makeatletter
% figure registry (figdefs.tex is generated by the extract step)
\newcommand{\qfigdef}[4]{\@namedef{qf@file@#1}{#2}\@namedef{qf@w@#1}{#3}\@namedef{qf@wide@#1}{#4}}
\newcommand{\qf@check}[1]{\@ifundefined{qf@file@#1}{\PackageError{score2}{Unknown figure id #1}{}}{}}
\newcommand{\qf@img}[4][]{\includegraphics[width=#2\dimexpr\@nameuse{qf@w@#4}pt\relax,max width=#3,max height=0.45\textheight,#1]{\@nameuse{qf@file@#4}}}
% Full-width questions: finalize wraps every question that contains a wide figure in \qfullbegin...\qfullend,
% so the WHOLE question (stem, figures, options) is set across the page in reading order, then the two columns
% resume. The question's head - its number, stem and first figure (\qfullsplit marks the end of the head) -
% is kept on one page, so a number is never left at a page foot with its figure overleaf; the rest may break
% across pages normally, so no big gaps are left. Without a \qfullsplit, the whole question is kept together
% if it is under 45% of a page. Never write these yourself.
\newif\ifqf@full
\newif\ifqf@boxed
\newsavebox{\qf@box}
\newcommand{\qf@open}{\global\qf@fulltrue\par\global\qf@boxedtrue
  \setbox\qf@box\vbox\bgroup\hsize\textwidth\linewidth\textwidth}
\newcommand{\qf@emit}[1]{\par\egroup\global\qf@boxedfalse
  \ifdim\dimexpr\ht\qf@box+\dp\qf@box\relax<#1\textheight\noindent\usebox{\qf@box}\par\else\unvbox\qf@box\fi}
\newcommand{\qfullsplit}{\ifqf@boxed\qf@emit{0.85}\fi}
\newcommand{\qf@close}{\ifqf@boxed\qf@emit{0.45}\fi\global\qf@fullfalse}
\newcommand{\qfullbegin}{\end{multicols}\qf@open}
\newcommand{\qfullend}{\qf@close\begin{multicols}{2}}
\newcommand{\qfullstart}{\qf@open}
\newcommand{\qfulldone}{\qf@close}
% \qfig[scale]{id}: block figure, centred, directly where it is written (size set by the tool).
% Fallback if a wide figure is not inside a full-width question: break out of the columns in place - no floats.
\newcommand{\qfig}[2][1]{\qf@check{#2}%
  \ifqf@full\qf@blk{#1}{#2}\else
  \ifnum\@nameuse{qf@wide@#2}=1 \expandafter\@firstoftwo\else\expandafter\@secondoftwo\fi
  {\end{multicols}\begin{center}\qf@img{#1}{\textwidth}{#2}\end{center}\begin{multicols}{2}}%
  {\qf@blk{#1}{#2}}\fi}
\newcommand{\qf@blk}[2]{\par\smallskip{\centering\qf@img{#1}{\linewidth}{#2}\par}\smallskip}
% \qfigi[scale]{id}: inline figure (inside option lists / tables / minipages); never breaks columns.
\newcommand{\qfigi}[2][1]{\qf@check{#2}\qf@img[valign=c]{#1}{\linewidth}{#2}}
\makeatother
% structure macros - lint reads these to check completeness and figure placement
\newcommand{\papersec}[1]{\par\bigskip{\centering\Large\bfseries PAPER-#1\par}\medskip}
\newcommand{\partsec}[2]{\par\medskip\noindent{\large\bfseries PART-#1: #2}\par\smallskip}
\newcommand{\sect}[1]{\par\smallskip\noindent\textbf{#1}\par}
\newcommand{\qn}[1]{\par\smallskip\noindent\textbf{#1.}~}
\newcommand{\paraq}[2]{\par\smallskip\noindent\textbf{Paragraph for Questions #1 and #2}\par}
\newenvironment{qoptions}{\begin{itemize}[leftmargin=1.8em,itemsep=1pt,topsep=2pt,label={}]}{\end{itemize}}
\newcommand{\opt}[2]{\item[(#1)] #2}
% figure registry (generated by score2_tools.py finalize)
\qfigdef{p1_phy_s1_q01_1}{TOAS - Test 15 (22 Aug 2022)-img/p1_phy_s1_q01_1.png}{44.9}{0}
\qfigdef{p1_phy_s1_q02_1}{TOAS - Test 15 (22 Aug 2022)-img/p1_phy_s1_q02_1.png}{132.2}{0}
\qfigdef{p1_phy_s1_q03_1}{TOAS - Test 15 (22 Aug 2022)-img/p1_phy_s1_q03_1.png}{103.2}{0}
\qfigdef{p1_phy_s1_q04_1}{TOAS - Test 15 (22 Aug 2022)-img/p1_phy_s1_q04_1.png}{157.9}{0}
\qfigdef{p1_phy_s1_q05_1}{TOAS - Test 15 (22 Aug 2022)-img/p1_phy_s1_q05_1.png}{89.9}{0}
\qfigdef{p1_phy_s2_q07_1}{TOAS - Test 15 (22 Aug 2022)-img/p1_phy_s2_q07_1.png}{84.6}{0}
\qfigdef{p1_phy_s2_q08_1}{TOAS - Test 15 (22 Aug 2022)-img/p1_phy_s2_q08_1.png}{166.1}{0}
\qfigdef{p1_phy_s2_q09_1}{TOAS - Test 15 (22 Aug 2022)-img/p1_phy_s2_q09_1.png}{132.2}{0}
\qfigdef{p1_phy_s2_q10_1}{TOAS - Test 15 (22 Aug 2022)-img/p1_phy_s2_q10_1.png}{179.6}{0}
\qfigdef{p1_phy_s2_q11_1}{TOAS - Test 15 (22 Aug 2022)-img/p1_phy_s2_q11_1.png}{100.1}{0}
\qfigdef{p1_phy_s2_q12_1}{TOAS - Test 15 (22 Aug 2022)-img/p1_phy_s2_q12_1.png}{150.3}{0}
\qfigdef{p1_phy_s3_q05_1}{TOAS - Test 15 (22 Aug 2022)-img/p1_phy_s3_q05_1.png}{84.4}{0}
\qfigdef{p1_phy_s3_q06_1}{TOAS - Test 15 (22 Aug 2022)-img/p1_phy_s3_q06_1.png}{103.2}{0}
\qfigdef{p1_chem_s1_q01_1}{TOAS - Test 15 (22 Aug 2022)-img/p1_chem_s1_q01_1.png}{161.3}{0}
\qfigdef{p1_chem_s1_q01_A}{TOAS - Test 15 (22 Aug 2022)-img/p1_chem_s1_q01_A.png}{92.6}{0}
\qfigdef{p1_chem_s1_q01_B}{TOAS - Test 15 (22 Aug 2022)-img/p1_chem_s1_q01_B.png}{108.3}{0}
\qfigdef{p1_chem_s1_q01_C}{TOAS - Test 15 (22 Aug 2022)-img/p1_chem_s1_q01_C.png}{108.5}{0}
\qfigdef{p1_chem_s1_q01_D}{TOAS - Test 15 (22 Aug 2022)-img/p1_chem_s1_q01_D.png}{108.0}{0}
\qfigdef{p1_chem_s1_q02_1}{TOAS - Test 15 (22 Aug 2022)-img/p1_chem_s1_q02_1.png}{170.2}{0}
\qfigdef{p1_chem_s1_q02_A}{TOAS - Test 15 (22 Aug 2022)-img/p1_chem_s1_q02_A.png}{74.2}{0}
\qfigdef{p1_chem_s1_q02_B}{TOAS - Test 15 (22 Aug 2022)-img/p1_chem_s1_q02_B.png}{74.2}{0}
\qfigdef{p1_chem_s1_q02_C}{TOAS - Test 15 (22 Aug 2022)-img/p1_chem_s1_q02_C.png}{79.4}{0}
\qfigdef{p1_chem_s1_q02_D}{TOAS - Test 15 (22 Aug 2022)-img/p1_chem_s1_q02_D.png}{68.9}{0}
\qfigdef{p1_chem_s1_q03_1}{TOAS - Test 15 (22 Aug 2022)-img/p1_chem_s1_q03_1.png}{187.6}{0}
\qfigdef{p1_chem_s1_q03_A}{TOAS - Test 15 (22 Aug 2022)-img/p1_chem_s1_q03_A.png}{48.9}{0}
\qfigdef{p1_chem_s1_q03_B}{TOAS - Test 15 (22 Aug 2022)-img/p1_chem_s1_q03_B.png}{53.9}{0}
\qfigdef{p1_chem_s1_q03_C}{TOAS - Test 15 (22 Aug 2022)-img/p1_chem_s1_q03_C.png}{52.9}{0}
\qfigdef{p1_chem_s1_q03_D}{TOAS - Test 15 (22 Aug 2022)-img/p1_chem_s1_q03_D.png}{48.2}{0}
\qfigdef{p1_chem_s2_q10_A}{TOAS - Test 15 (22 Aug 2022)-img/p1_chem_s2_q10_A.png}{77.6}{0}
\qfigdef{p1_chem_s2_q10_B}{TOAS - Test 15 (22 Aug 2022)-img/p1_chem_s2_q10_B.png}{110.0}{0}
\qfigdef{p1_chem_s2_q10_C}{TOAS - Test 15 (22 Aug 2022)-img/p1_chem_s2_q10_C.png}{131.3}{0}
\qfigdef{p1_chem_s2_q11_1}{TOAS - Test 15 (22 Aug 2022)-img/p1_chem_s2_q11_1.png}{33.5}{0}
\qfigdef{p1_chem_s2_q11_2}{TOAS - Test 15 (22 Aug 2022)-img/p1_chem_s2_q11_2.png}{50.2}{0}
\qfigdef{p1_chem_s2_q11_A}{TOAS - Test 15 (22 Aug 2022)-img/p1_chem_s2_q11_A.png}{71.3}{0}
\qfigdef{p1_chem_s2_q11_B}{TOAS - Test 15 (22 Aug 2022)-img/p1_chem_s2_q11_B.png}{57.8}{0}
\qfigdef{p1_chem_s2_q11_C}{TOAS - Test 15 (22 Aug 2022)-img/p1_chem_s2_q11_C.png}{66.1}{0}
\qfigdef{p1_chem_s2_q11_D}{TOAS - Test 15 (22 Aug 2022)-img/p1_chem_s2_q11_D.png}{28.6}{0}
\qfigdef{p1_chem_s2_q11_7D}{TOAS - Test 15 (22 Aug 2022)-img/p1_chem_s2_q11_7D.png}{28.9}{0}
\qfigdef{p1_chem_s2_q12_A}{TOAS - Test 15 (22 Aug 2022)-img/p1_chem_s2_q12_A.png}{59.0}{0}
\qfigdef{p1_chem_s2_q12_B}{TOAS - Test 15 (22 Aug 2022)-img/p1_chem_s2_q12_B.png}{58.5}{0}
\qfigdef{p1_chem_s2_q12_C}{TOAS - Test 15 (22 Aug 2022)-img/p1_chem_s2_q12_C.png}{61.7}{0}
\qfigdef{p1_chem_s2_q12_D}{TOAS - Test 15 (22 Aug 2022)-img/p1_chem_s2_q12_D.png}{67.6}{0}
\qfigdef{p1_chem_s3_q02_1}{TOAS - Test 15 (22 Aug 2022)-img/p1_chem_s3_q02_1.png}{145.4}{0}
\qfigdef{p1_chem_s3_q03_1}{TOAS - Test 15 (22 Aug 2022)-img/p1_chem_s3_q03_1.png}{183.7}{0}
\qfigdef{p2_phy_s1_q01_1}{TOAS - Test 15 (22 Aug 2022)-img/p2_phy_s1_q01_1.png}{82.8}{0}
\qfigdef{p2_phy_s1_q03_1}{TOAS - Test 15 (22 Aug 2022)-img/p2_phy_s1_q03_1.png}{82.5}{0}
\qfigdef{p2_phy_s1_q04_1}{TOAS - Test 15 (22 Aug 2022)-img/p2_phy_s1_q04_1.png}{158.7}{0}
\qfigdef{p2_phy_s1_q05_1}{TOAS - Test 15 (22 Aug 2022)-img/p2_phy_s1_q05_1.png}{128.6}{0}
\qfigdef{p2_phy_s1_q06_1}{TOAS - Test 15 (22 Aug 2022)-img/p2_phy_s1_q06_1.png}{96.7}{0}
\qfigdef{p2_phy_s2_q04_1}{TOAS - Test 15 (22 Aug 2022)-img/p2_phy_s2_q04_1.png}{105.3}{0}
\qfigdef{p2_phy_s2_q05_1}{TOAS - Test 15 (22 Aug 2022)-img/p2_phy_s2_q05_1.png}{99.7}{0}
\qfigdef{p2_phy_s2_q06_1}{TOAS - Test 15 (22 Aug 2022)-img/p2_phy_s2_q06_1.png}{106.0}{0}
\qfigdef{p2_phy_s3_q01_1}{TOAS - Test 15 (22 Aug 2022)-img/p2_phy_s3_q01_1.png}{84.2}{0}
\qfigdef{p2_phy_s3_q02_1}{TOAS - Test 15 (22 Aug 2022)-img/p2_phy_s3_q02_1.png}{156.5}{0}
\qfigdef{p2_phy_s3_q03_1}{TOAS - Test 15 (22 Aug 2022)-img/p2_phy_s3_q03_1.png}{84.4}{0}
\qfigdef{p2_phy_s3_q04_1}{TOAS - Test 15 (22 Aug 2022)-img/p2_phy_s3_q04_1.png}{99.9}{0}
\qfigdef{p2_phy_s3_q06_1}{TOAS - Test 15 (22 Aug 2022)-img/p2_phy_s3_q06_1.png}{121.8}{0}
\qfigdef{p2_chem_s1_q01_A}{TOAS - Test 15 (22 Aug 2022)-img/p2_chem_s1_q01_A.png}{64.1}{0}
\qfigdef{p2_chem_s1_q01_B}{TOAS - Test 15 (22 Aug 2022)-img/p2_chem_s1_q01_B.png}{40.4}{0}
\qfigdef{p2_chem_s1_q01_C}{TOAS - Test 15 (22 Aug 2022)-img/p2_chem_s1_q01_C.png}{47.7}{0}
\qfigdef{p2_chem_s1_q01_D}{TOAS - Test 15 (22 Aug 2022)-img/p2_chem_s1_q01_D.png}{39.7}{0}
\qfigdef{p2_chem_s1_q02_A}{TOAS - Test 15 (22 Aug 2022)-img/p2_chem_s1_q02_A.png}{84.5}{0}
\qfigdef{p2_chem_s1_q02_C}{TOAS - Test 15 (22 Aug 2022)-img/p2_chem_s1_q02_C.png}{122.5}{0}
\qfigdef{p2_chem_s1_q02_D}{TOAS - Test 15 (22 Aug 2022)-img/p2_chem_s1_q02_D.png}{137.0}{0}
\qfigdef{p2_chem_s1_q03_A}{TOAS - Test 15 (22 Aug 2022)-img/p2_chem_s1_q03_A.png}{29.2}{0}
\qfigdef{p2_chem_s1_q03_2}{TOAS - Test 15 (22 Aug 2022)-img/p2_chem_s1_q03_2.png}{32.0}{0}
\qfigdef{p2_chem_s1_q03_B}{TOAS - Test 15 (22 Aug 2022)-img/p2_chem_s1_q03_B.png}{87.2}{0}
\qfigdef{p2_chem_s1_q03_C}{TOAS - Test 15 (22 Aug 2022)-img/p2_chem_s1_q03_C.png}{84.6}{0}
\qfigdef{p2_chem_s1_q05_1}{TOAS - Test 15 (22 Aug 2022)-img/p2_chem_s1_q05_1.png}{187.1}{0}
\qfigdef{p2_chem_s2_q03_1}{TOAS - Test 15 (22 Aug 2022)-img/p2_chem_s2_q03_1.png}{196.4}{0}
\qfigdef{p2_chem_s2_q03_2}{TOAS - Test 15 (22 Aug 2022)-img/p2_chem_s2_q03_2.png}{192.5}{0}
\qfigdef{p2_chem_s2_q03_3}{TOAS - Test 15 (22 Aug 2022)-img/p2_chem_s2_q03_3.png}{177.5}{0}
\qfigdef{p2_chem_s2_q04_1}{TOAS - Test 15 (22 Aug 2022)-img/p2_chem_s2_q04_1.png}{175.4}{0}
\qfigdef{p2_chem_s2_q05_2}{TOAS - Test 15 (22 Aug 2022)-img/p2_chem_s2_q05_2.png}{56.3}{0}
\qfigdef{p2_chem_s2_q05_1}{TOAS - Test 15 (22 Aug 2022)-img/p2_chem_s2_q05_1.png}{27.9}{0}
\qfigdef{p2_chem_s2_q05_3}{TOAS - Test 15 (22 Aug 2022)-img/p2_chem_s2_q05_3.png}{43.2}{0}
\qfigdef{p2_chem_s2_q05_5}{TOAS - Test 15 (22 Aug 2022)-img/p2_chem_s2_q05_5.png}{84.5}{0}
\qfigdef{p2_chem_s2_q05_4}{TOAS - Test 15 (22 Aug 2022)-img/p2_chem_s2_q05_4.png}{48.7}{0}
\qfigdef{p2_chem_s2_q05_7}{TOAS - Test 15 (22 Aug 2022)-img/p2_chem_s2_q05_7.png}{58.8}{0}
\qfigdef{p2_chem_s2_q05_6}{TOAS - Test 15 (22 Aug 2022)-img/p2_chem_s2_q05_6.png}{51.9}{0}
\raggedcolumns
\makeatletter\@beginparpenalty=10000\makeatother
\newcolumntype{L}[1]{>{\raggedright\arraybackslash}p{#1}}
\begin{document}
\begin{center}{\Large\bfseries TOAS --- Test 15 (22 Aug 2022)}\end{center}
\begin{multicols}{2}
\raggedcolumns
\papersec{1}
\partsec{1}{PHYSICS}
\sect{SECTION-I(i)}
\qn{1} A long conducting wire is placed in plane of a ring of radius $R$ at distance $d\approx R$ from centre. Current in ring $i$ is increasing at constant rate $\alpha$. Then induced emf in wire :
\qfig{p1_phy_s1_q01_1}
\begin{qoptions}
\opt{A}{$\mu_0R\alpha$}
\opt{B}{$\dfrac{\mu_0R\alpha}{2}$}
\opt{C}{$\pi\mu_0R\alpha$}
\opt{D}{$\dfrac{\mu_0R\alpha}{\pi}$}
\end{qoptions}
\qn{2} A conducting wire of (mass 0.1 kg) is placed on conducting rails inclined at angle $37^\circ$ from horizontal. Ends of rails are connected by 5V, $5\Omega$ battery when magnetic field $B=2$T is switch on in vertically upward direction then acceleration of wire just after magnetic field is switch on \{self inductance of current loop is 0.2 H initially, friction coefficient between wire and rails is 3/4\}.
\qfig{p1_phy_s1_q02_1}
\begin{qoptions}
\opt{A}{0}
\opt{B}{2.95 m/s$^2$}
\opt{C}{1.25 m/s$^2$}
\opt{D}{4/15 m/s$^2$}
\end{qoptions}
\qn{3} Reading of voltmeter and ammeter is :-
\qfig{p1_phy_s1_q03_1}
\begin{qoptions}
\opt{A}{$10\sqrt2$ volt, $\sqrt{10}$A}
\opt{B}{10 volt, $\sqrt5$A}
\opt{C}{$10\sqrt2$ volt, $\sqrt5$A}
\opt{D}{10 volt, $\sqrt{10}$A}
\end{qoptions}
\qn{4} Value of current $I$ is :-
\qfig{p1_phy_s1_q04_1}
\begin{qoptions}
\opt{A}{5A}
\opt{B}{$\left(\dfrac{15}{4}\right)$A}
\opt{C}{$\left(\dfrac{45}{64}\right)$A}
\opt{D}{None}
\end{qoptions}
\qn{5} A small ball of mass $m$, charge $Q$ is to be held in stable equilibrium for small angular displacement as shown in figure. Find minimum charge $q$ that should be fixed below ball for this :-
\qfig{p1_phy_s1_q05_1}
\begin{qoptions}
\opt{A}{$\dfrac{16\pi\epsilon_0mgR^2}{Q}$}
\opt{B}{$\dfrac{32\pi\epsilon_0mgR^2}{Q}$}
\opt{C}{$\dfrac{8\pi\epsilon_0mgR^2}{3Q}$}
\opt{D}{$\dfrac{64\pi\epsilon_0mgR^2}{Q}$}
\end{qoptions}
\qn{6} An asteroid of mass $M$ exploded into a spherical homogeneous cloud [spherical shell] in free space. Due to energy received by the explosion, the cloud expands and the expansion is spherically symmetric. At an instant, when radius of the cloud is $R_0$, all of its particles on the surface are observed receding radially away from the centre of the cloud with a velocity $v_0$. What will be the radius of the cloud, when its expansion ceases?
\begin{qoptions}
\opt{A}{$\dfrac{GMR_0}{GM-v_0^2R_0}$}
\opt{B}{$\dfrac{GM}{v_0^2}$}
\opt{C}{$\dfrac{2GMR_0}{2GM-v_0^2R_0}$}
\opt{D}{$\dfrac{GMR_0}{2GM-V_0^2R_0}$}
\end{qoptions}
\sect{SECTION-I(ii)}
\qn{7} Three concentric conducting spherical shells have radius $R$, $2R$ \& $3R$ are shown. Middle shell is earthed and inner most shell is connected to outer most shell by conducting wire. Space between $2^{\text{nd}}$ and $3^{\text{rd}}$ shell is filled by dielectric of dielectric constant $K=2$ and $Q=17\mu$C charge is given to $3^{\text{rd}}$ shell then :-
\qfig{p1_phy_s2_q07_1}
\begin{qoptions}
\opt{A}{Capacitance of system is $44\pi\epsilon_0R$}
\opt{B}{Capacitance of system is $68\pi\epsilon_0R$}
\opt{C}{Charge on $2^{\text{nd}}$ shell is $-14\mu$C.}
\opt{D}{Charge on $2^{\text{nd}}$ shell is $-17\mu$C.}
\end{qoptions}
\qn{8} The figure shows a snap shot of a vibrating string of linear mass density 0.1 kg/m at $t=0$. The particle P is observed moving up with velocity $20\sqrt3$cm/s. The tangent at P makes an angle $60^\circ$ with x-axis then
\qfig{p1_phy_s2_q08_1}
\begin{qoptions}
\opt{A}{The wave equation is given by $y=4\text{mm}\sin\left[10\pi t+50\pi x+\dfrac\pi4\right]$.}
\opt{B}{The wave equation is given by $y=4\text{mm}\cos\left[10\pi t-50\pi x+\dfrac\pi4\right]$.}
\opt{C}{Rate of flow of energy is towards left and it's value at point P is $12\times10^{-5}$J}
\opt{D}{Tension in string is $4\times10^{-3}$ N.}
\end{qoptions}
\qn{9} Two string having linear mass density 0.1 kg/m and 0.4kg/m are connected as shown and pulled by force 10 N. Select correct alternative :-
\qfig{p1_phy_s2_q09_1}
\begin{qoptions}
\opt{A}{Fundamental frequency of system is 10Hz, if antinote is formed at point P.}
\opt{B}{Fundamental frequency of system is 5Hz, if node is formed at point P.}
\opt{C}{$1^{\text{st}}$ overtone has frequency 30Hz, if antinote is formed at point P.}
\opt{D}{$1^{\text{st}}$ overtone has frequency 10Hz, if node is formed at point P.}
\end{qoptions}
\qn{10} In given situation find possible values of distance $x$ between lens and mirror so that final image forms at object itself. Refractive index of water is $4/3$.
\qfig{p1_phy_s2_q10_1}
\begin{qoptions}
\opt{A}{96 cm}
\opt{B}{36 cm}
\opt{C}{66 cm}
\opt{D}{156 cm}
\end{qoptions}
\qn{11} A wire having a uniform linear charge density $\lambda$, is bent in the form of ring of radius $R$. Point A as shown in the figure, is in the plane of ring but not at centre. Two elements of the ring of length $a_1$ and $a_2$ subtend very small same angle at the point A. They are at distance $r_1$ and $r_2$ from the point A respectively.
\qfig{p1_phy_s2_q11_1}
\begin{qoptions}
\opt{A}{The ratio of charge of elements $a_1$ and $a_2$ is $\dfrac{r_1}{r_2}$.}
\opt{B}{The element $a_1$ produced greater magnitude of electric field at A than $a_2$.}
\opt{C}{The element $a_1$ and $a_2$ produced same potential at A.}
\opt{D}{The direction of net electric field at A is towards element $a_2$.}
\end{qoptions}
\qn{12} In given figure the wires $P_1Q_1$ and $P_2Q_2$ are made to slide on the rails with same speed 5 cm/s. In this region a uniform magnetic field 1T exists inside plane of paper. The current in $9\Omega$ resistor is :
\qfig{p1_phy_s2_q12_1}
\begin{qoptions}
\opt{A}{Zero if both wires slide towards left.}
\opt{B}{Zero if both wires slides in opposite direction.}
\opt{C}{0.2 mA if both wire move towards left.}
\opt{D}{0.2 mA if both wire move in opposite direction.}
\end{qoptions}
\sect{SECTION-II}
\qn{1} A water jet shoot water with speed 5 m/s from an opening of area 15 cm$^2$ in upward direction. A plate is balanced at height 80 cm above jet by force applied by water. After colliding the plate water goes horizontally in every possible direction. Assuming water to be ideal fluid find angular frequency (in rad/s) of small oscillation of plate in vertical direction. (Take $g=10$ m/s$^2$, density of water $=10^3$ kg/m$^3$)
\qn{2} A current $I=10$A flows in a long thin walled cylinder of radius 1cm. Then pressure on walls due to magnetic force is (in Pa)
\qn{3} A particle is projected from point, that is at a distance $4R$ from centre of the earth, with speed $50\sqrt2\,x$ m/s, in a direction making angle $30^\circ$ with vertical, away from the earth. Find $x$ if particle passes grazing the surface of the earth use $\dfrac{GM}{R}=6.4\times10^7$m$^2$/s$^2$, where $R$ is radius of earth and $M$ is mass of the earth.
\qn{4} Position of a particle moving along x axis is given by $x=A\cos\left(\dfrac{\pi}{180}t\right)$. Find time (in sec) in which particle will travel a distance $\dfrac{12A}{5}$.
\qn{5} A transparent hexagonal body ABCDEF is shown find minimum value of refractive index of body so that ray incidenting at surface AB do not reach observer through surface CD.
\qfig{p1_phy_s3_q05_1}
\qn{6} At $t=0$, key S is closed. If $I_1$ represent current through cell just after key is closed while $I_2$ represent current long after key is closed then $I_1+I_2$ will be (in A).
\qfig{p1_phy_s3_q06_1}
\partsec{2}{CHEMISTRY}
\sect{SECTION-I(i)}
\qn{1}
\qfig{p1_chem_s1_q01_1}
in above reaction Y, A, B respectively is :-
\begin{qoptions}
\opt{A}{\qfigi{p1_chem_s1_q01_A}}
\opt{B}{\qfigi{p1_chem_s1_q01_B}}
\opt{C}{\qfigi{p1_chem_s1_q01_C}}
\opt{D}{\qfigi{p1_chem_s1_q01_D}}
\end{qoptions}
\qn{2}
\qfig{p1_chem_s1_q02_1}
\begin{qoptions}
\opt{A}{\qfigi{p1_chem_s1_q02_A}}
\opt{B}{\qfigi{p1_chem_s1_q02_B}}
\opt{C}{\qfigi{p1_chem_s1_q02_C}}
\opt{D}{\qfigi{p1_chem_s1_q02_D}}
\end{qoptions}
\qn{3}
\qfig{p1_chem_s1_q03_1}
\begin{qoptions}
\opt{A}{\qfigi{p1_chem_s1_q03_A}}
\opt{B}{\qfigi{p1_chem_s1_q03_B}}
\opt{C}{\qfigi{p1_chem_s1_q03_C}}
\opt{D}{\qfigi{p1_chem_s1_q03_D}}
\end{qoptions}
\qn{4} Which of the following is INCORRECT?
\begin{qoptions}
\opt{A}{H$_2$S does not give test with sodium nitroprusside}
\opt{B}{HgCl$_2$ does not give chromyl chloride}
\opt{C}{Mg$^{2+}$ can not be detected by flame test}
\opt{D}{Through blue glass brick red flame is observe for Ca$^{+2}$ in flame test}
\end{qoptions}
\qn{5} Which of the following set is correct?
\begin{qoptions}
\opt{A}{L.D. Process $\to$ Used for : Manufacturing of steel}
\opt{B}{Hall's Heroult Process $\to$ Used for : Conversion of bauxite ore into pure alumina}
\opt{C}{Poling Process $\to$ Used for : Purification of impure Cu}
\opt{D}{Liquation Process $\to$ Used for : Purification of impure Pb}
\end{qoptions}
\qn{6} Which of the following having density less than Na?
\begin{qoptions}
\opt{A}{K}
\opt{B}{Rb}
\opt{C}{Cs}
\opt{D}{None of these}
\end{qoptions}
\sect{SECTION-I(ii)}
\qn{7} The standard boiling point of A($\ell$) is 400K. The enthalpy of vapourisation A($\ell$) at 400 K and 1 bar is 20 kcal/ mol. Which of the following information (s) is/are correct regarding the process\par
\noindent A($\ell$, 400K, 1 bar) $\longrightarrow$ A(g, 400K, 0.5 bar)\par
[Given : $C_{p,m}$ values (in cal / K-mol) : A(g) = 5 ; A(l) = 10; $\ln2=0.7$, $\ln5=1.6$ , Assume ideal behaviour of A(g)]
\begin{qoptions}
\opt{A}{$\Delta H_m=+20$ kcal/ mol}
\opt{B}{$\Delta S_m=+45.1$ cal/ K-mol}
\opt{C}{$\Delta G_m=-560$ cal/ mol}
\opt{D}{$\Delta U_m=20$ kcal/mole}
\end{qoptions}
\qn{8} Select the correct statements
\begin{qoptions}
\opt{A}{Catalyst can change the percentage yield of a reversible reaction}
\opt{B}{Molecularity and order are same for the elementary reactions}
\opt{C}{The catalyst forms an activated complex of lower potential energy, with the reactants by which more number of molecules are able to cross the barrier per unit of time.}
\opt{D}{For exothermic reaction rate constant decreases with increasing temperature}
\end{qoptions}
\qn{9} Consider the following reaction :\par
$2\text{Fe}_2\text{O}_3(s)+3\text{C}(s)\to4\text{Fe}(s)+3\text{CO}_2(g)$\par
$\Delta H_f^\circ$ of Fe$_2$O$_3$ \& CO$_2$ are $-820$ kJ/mol and $-390$ kJ/mol respectively.\par
The reaction is -
\begin{qoptions}
\opt{A}{Endothermic}
\opt{B}{Exothermic}
\opt{C}{Spontaneous at high temperature}
\opt{D}{Spontaneous at low temperature}
\end{qoptions}
\qn{10} $\text{P}+\text{HCHO}\xrightarrow{\text{OH}^-/\text{ Polymerisation}}$ Bakelite\par
Find the reaction in which P form as product.
\begin{qoptions}
\opt{A}{\qfigi{p1_chem_s2_q10_A}}
\opt{B}{\qfigi{p1_chem_s2_q10_B}}
\opt{C}{\qfigi{p1_chem_s2_q10_C}}
\opt{D}{None}
\end{qoptions}
\qn{11} (P) \qfigi{p1_chem_s2_q11_1}\qfigi{p1_chem_s2_q11_2}\par
Possible P and Q is/are :-
\begin{qoptions}
\opt{A}{\qfigi{p1_chem_s2_q11_A}}
\opt{B}{\qfigi{p1_chem_s2_q11_B}}
\opt{C}{\qfigi{p1_chem_s2_q11_C}}
\opt{D}{\qfigi{p1_chem_s2_q11_D}\qfigi{p1_chem_s2_q11_7D}}
\end{qoptions}
\qn{12} An ideal gas undergoes a reversible isothermal expansion from state I to state II followed by a reversible adiabatic expansion from state II to state III. The correct plot(s) representing the changes from state I to state III is(are)\par
(p: pressure, V: volume, T: temperature, H: enthalpy, S: entropy)
\begin{qoptions}
\opt{A}{\qfigi{p1_chem_s2_q12_A}}
\opt{B}{\qfigi{p1_chem_s2_q12_B}}
\opt{C}{\qfigi{p1_chem_s2_q12_C}}
\opt{D}{\qfigi{p1_chem_s2_q12_D}}
\end{qoptions}
\sect{SECTION-II}
\qn{1} 1 molal aqueous solution of an electrolyte AB$_2$ is 80\% ionised in to A$^{+2}$ and B$^{-1}$ ions. The boiling point of the solution (in Kelvin) at 1 atm is : ($K_b(\text{H}_2\text{O})=0.5$ K kg mol$^{-1}$, 0 K $=-273^\circ$C).
\qn{2} For the dissociation of SO$_3$ into SO$_2$ and O$_2$, at a given temperature, the equilibrium constant $K_p=x$ atm$^2$, then find $x$. The variation of the partial pressure of the three gases with time is as shown :
\qfig{p1_chem_s3_q02_1}
\qn{3} Examine the structure of following compounds, and find out number of compounds that will undergoes decarboxylation in presence of heat.
\qfig{p1_chem_s3_q03_1}
\qn{4} Calculate the magnetic moment of $[\text{Co}(\text{H}_2\text{O})_6]^{3+}$ in B.M.
\qn{5} What is the average oxidation of `S' in product in I$_2$ reacts with hypo solution.
\qn{6} The crystal field stabilizing energy of $[\text{IrF}_6]^{3-}$ is $-x\Delta_0+y\text{P}$.\par
find the sum of the magnitude of $x+y$ is
\partsec{3}{MATHEMATICS}
\sect{SECTION-I(i)}
\qn{1} Let $A=\begin{bmatrix}x^2\\1\\2x\end{bmatrix}$, $B=[\,a\ \ b\ \ c\,]$ and $C=\begin{bmatrix}3x^2&5x^2&2x\\2x&4&3x\\5x&7&6x\end{bmatrix}$, given $\text{tr}(AB)=\text{tr}(C)\ \forall\ x\in R$, where $\text{tr}(A)$ denotes trace of $A$, then the value of $\displaystyle\int_0^\infty\dfrac{\ell nx}{ax^2+bx+c}\,dx$
\begin{qoptions}
\opt{A}{0}
\opt{B}{1}
\opt{C}{$-3$}
\opt{D}{3}
\end{qoptions}
\qn{2} The integral curve through origin satisfying the differential equation $y^2(3x^9+y)\dfrac{dy}{dx}=x^3(1-9x^5y^3)$ is
\begin{qoptions}
\opt{A}{$3(x^3y)^2+y^3-x^3=0$}
\opt{B}{$4(x^3y)^3+y^4-x^4=0$}
\opt{C}{$4(x^3y)^3-y^4+x^4=0$}
\opt{D}{$3(x^3y)^2-y^3+x^3=0$}
\end{qoptions}
\qn{3} If $A=\begin{bmatrix}1&0&0\\0&1&2\\0&-2&3\end{bmatrix}$ \& $I=\begin{bmatrix}1&0&0\\0&1&0\\0&0&1\end{bmatrix}$, then if $A^7-4A^6+6A^5=\alpha A^2+\beta A+\gamma I$ then $(\alpha,\beta,\gamma)$ is :
\begin{qoptions}
\opt{A}{$(5,-11,7)$}
\opt{B}{$(-21,115,-91)$}
\opt{C}{$(-13,44,-28)$}
\opt{D}{None of these}
\end{qoptions}
\qn{4} If the solution of the differential equation $\dfrac{x\,dx-y\,dy}{x\,dy-y\,dx}=\sqrt{\dfrac{1+x^2-y^2}{x^2-y^2}}$ be $\sqrt{f(x,y)}+\sqrt{1+f(x,y)}=c\left(\dfrac{x+y}{\sqrt{f(x,y)}}\right)$, then $f(x,y)$ is :
\begin{qoptions}
\opt{A}{$x^2+y^2$}
\opt{B}{$1+x^2-y^2$}
\opt{C}{$x^2-y^2$}
\opt{D}{$\dfrac{x^2-y^2}{x^2+y^2}$}
\end{qoptions}
\qn{5} $\displaystyle\lim_{n\to\infty}\left(\sqrt{n^2+n+1}-\left[\sqrt{n^2+n+1}\right]\right)$ $(n\in N)$ (where $[.]$ denotes the greatest integer function) is
\begin{qoptions}
\opt{A}{0}
\opt{B}{1/2}
\opt{C}{2/3}
\opt{D}{1/4}
\end{qoptions}
\qn{6} A rack has 5 different pairs of shoes. The number of ways in which 4 shoes can be chosen from it, such that there will be no complete pair is :
\begin{qoptions}
\opt{A}{1920}
\opt{B}{200}
\opt{C}{110}
\opt{D}{80}
\end{qoptions}
\sect{SECTION-I(ii)}
\qn{7} If the solutions of the equation $x^3+1=2\sqrt[3]{2x-1}$ are $\left\{a,\dfrac{b\pm\sqrt c}{d}\right\}$, $(a,\ b\in I\ \&\ c,d$ are prime) then which of the following is(are) correct ?
\begin{qoptions}
\opt{A}{$a+b^3+c-d=3$}
\opt{B}{$a+b=1$}
\opt{C}{$c+d=7$}
\opt{D}{$c-d=3$}
\end{qoptions}
\qn{8} A function $f$ is defined by $f(x)=\displaystyle\int_0^\pi\cos t\,\cos(x-t)\,dt$, $0\le x\le2\pi$ then which of the following hold(s) good?
\begin{qoptions}
\opt{A}{$f(x)$ is continuous but not differentiable in $(0,2\pi)$}
\opt{B}{Maximum value of $f(x)$ is $\pi$}
\opt{C}{There exists atleast one $c\in(0,2\pi)$ such that $f'(c)=0$}
\opt{D}{Minimum value of $f(x)$ is $-\dfrac\pi2$}
\end{qoptions}
\qn{9} Given a function `$g$' which has a derivative $g'(0)=2$ and $g(x+y)=e^yg(x)+e^xg(y)$ for all $x$ and $y$, then:
\begin{qoptions}
\opt{A}{$g(x)$ is increasing for all $x\in[-1,\infty)$}
\opt{B}{Range of $g(x)$ is $\left[-\dfrac2e,\infty\right)$}
\opt{C}{$g''(x)>0\ \forall\ x\in R$}
\opt{D}{$\displaystyle\lim_{x\to0}\dfrac{g(x)}{x}=1$}
\end{qoptions}
\qn{10} If $y_1,y_2$ are the solutions of the differential equation $\dfrac{dy}{dx}+P(x)\,y=Q(x)$, then :
\begin{qoptions}
\opt{A}{$y=y_1+c(y_1-y_2)$ is the general solution of equation}
\opt{B}{$y=y_1+c(y_1+y_2)$ is the general solution of equation}
\opt{C}{$\alpha y_1+\beta y_2$ is a solution if $\alpha+\beta=1$}
\opt{D}{$\alpha y_1-\beta y_2$ is a solution if $\alpha-\beta=1$}
\end{qoptions}
\qn{11} Area of region bounded between the curves $y=\dfrac1{1+x^2}$ and $y=1-\dfrac1{1+x^2}$ between its points of intersection is
\begin{qoptions}
\opt{A}{$\pi-2$}
\opt{B}{an irrational number}
\opt{C}{an integer}
\opt{D}{4}
\end{qoptions}
\qn{12} The combinatorial coefficient $^{n-1}C_p$ denotes
\begin{qoptions}
\opt{A}{The number of ways in which $n$ things of which $p$ are alike and rest different can be arranged in a circle}
\opt{B}{The number of ways in which $p$ different things can be selected out of $n$ different thing if a particular thing is always excluded}
\opt{C}{Number of ways in which $n$ alike balls can be distributed in $p$ different boxes so that no box remains empty, and each box can hold any number of balls}
\opt{D}{The number of ways in which $(n-2)$ white balls and $p$ black balls can be arranged in a line if black balls are separated, balls are alike except for the colour}
\end{qoptions}
\sect{SECTION-II}
\qn{1} Let solution of equation $|x^2-3x-4|+|x^2+2x-3|=|5x+1|$ is $x\in[a,b]\cup[c,d]$, then value of $|a|+|b|+|c|+|d|$
\qn{2} Consider the matrix $A,B,C,D$ with order $2\times3$, $3\times4$, $4\times4$, $4\times2$ respectively.\par
Let $x=(\alpha AB\gamma C^2D)^3$, where $\alpha$ \& $\gamma$ are scalars. Let $|x|=k|ABC^2D|^3$ \& if $k=(\alpha\gamma)^\lambda$, then $\lambda$ is (where $|P|$ denotes determinant value of matrix $P$)
\qn{3} If $A=\left.\dfrac{d\left(\sin^{-1}\left(\frac{x}{\sqrt{4+x^2}}\right)\right)}{d\left(\tan^{-1}\frac x2\right)}\right|_{x=0}$ and $B=\left.\dfrac{d\left(\sin^{-1}x\right)}{d\left(\tan^{-1}\frac{\sqrt{1-x^2}}{x}\right)}\right|_{x=\frac12}$, then the value of $(2A+3B^2)$ is
\qn{4} If $I_1=\displaystyle\int_0^1\dfrac{dx}{e^x(1+x)}$ and $I_2=\displaystyle\int_0^{\pi/4}\dfrac{e^{\tan^2\theta}}{(2-\tan^2\theta)}\,\dfrac{\sin\theta\,d\theta}{\cos^3\theta}$, then $\dfrac{I_1}{I_2}$ is \underline{\hspace{2cm}}.
\qn{5} The area bounded by the curves $y=xe^x$, $y=xe^{-x}$ and the lines $x=0$ and $x=1$ is (in \textit{square units}) :
\qn{6} Let $A=[a_{ij}]_{3\times3}$, $B=[b_{ij}]_{3\times3}$ where $b_{ij}=3^{i-j}a_{ij}$, $C=[c_{ij}]_{3\times3}$, where $c_{ij}=4^{i-j}b_{ij}$ be any three matrices. If $|A|=4$, then $|B|-|C|$ ($|X|$ denotes determinant of matrix $X$)
\newpage
\papersec{2}
\partsec{1}{PHYSICS}
\sect{SECTION-I}
\qn{1} A small solid ball of density $\rho$ is held inside at point A a cubical container of side $L$, filled with an ideal liquid of density $4\rho$ as shown in the figure. Now, if the container starts moving with constant acceleration $a$ horizontally and the ball is released from point A simultaneously, then
\qfig{p2_phy_s1_q01_1}
\begin{qoptions}
\opt{A}{For ball to hit the top of container at end Q, $a=3g$}
\opt{B}{For ball to hit the top of container at end Q, $a=2g$}
\opt{C}{Ball hits the top of container at end Q after a time $t=\sqrt{\dfrac{L}{3g}}$}
\opt{D}{Ball hits the top of container at end Q after a time $t=\sqrt{\dfrac{2L}{3g}}$}
\end{qoptions}
\qn{2} $Y(x,t)=\dfrac{0.8}{\left[(4x+5t)^2+5\right]}$ represents a moving pulse where $x$ and $y$ are in metres and $t$ in second. Then :
\begin{qoptions}
\opt{A}{pulse is moving in positive x-direction.}
\opt{B}{in 2s it will travel a distance of 2.5 m.}
\opt{C}{maximum displacement of particles is 0.16 m.}
\opt{D}{it is a symmetric pulse at $t=0$.}
\end{qoptions}
\qn{3} Two concentric spherical shells masses $m_1$ and $m_2$ and radii $r_1$ and $r_2$. Then
\qfig{p2_phy_s1_q03_1}
\begin{qoptions}
\opt{A}{outer shell will have no contribution in gravitational field at point P.}
\opt{B}{force on P is directed towards O.}
\opt{C}{force on P is $\dfrac{Gm_1m_2}{r^2}$.}
\opt{D}{force on P is $\dfrac{Gm_1m_3}{r^2}$.}
\end{qoptions}
\qn{4} A part of the circuit is shown in the figure
\qfig{p2_phy_s1_q04_1}
\begin{qoptions}
\opt{A}{power dissipated in $3\Omega$ resistance is 27 watt}
\opt{B}{$V_C-V_B=20$ V}
\opt{C}{$V_C-V_D=20$ V}
\opt{D}{power dissipated in $6\Omega$ resistance is 36 watt}
\end{qoptions}
\qn{5} If the electric field $E$ is plotted, with distance from $q_1$ along the line joining of two, then ($E$ is positive along +ve x-axis) it looks as shown in figure. From the plot we can say that
\qfig{p2_phy_s1_q05_1}
\begin{qoptions}
\opt{A}{$q_1$ and $q_2$ both are negative}
\opt{B}{$\left|\dfrac{q_1}{q_2}\right|=4$}
\opt{C}{$q_1$ is positive and $q_2$ is negative}
\opt{D}{$\left|\dfrac{q_1}{q_2}\right|=\dfrac14$}
\end{qoptions}
\qn{6} A container is partially filled with a liquid of density $\rho_2$. A capillary tube of radius `$r$' is vertically inserted in this liquid. Now another liquid of density $\rho_1(\rho_1<\rho_2)$ is slowly poured in the container to a height $h$ as shown. There is only denser liquid in the capillary tube. The rise of denser liquid in the capillary tube is also $h$. (Assume zero contact angle).
\qfig{p2_phy_s1_q06_1}
\begin{qoptions}
\opt{A}{the surface tension of heavier liquid is $2\pi r(\rho_2-\rho_1)gh$}
\opt{B}{the surface tension of heavier liquid is $\dfrac r2(\rho_2-\rho_1)gh$}
\opt{C}{liquid level in capillary will be lower than free surface of lighter liquid if surface tension is zero.}
\opt{D}{liquid level in capillary will be equal to the free surface of lighter liquid if surface tension is zero.}
\end{qoptions}
\sect{SECTION-II}
\qn{1} A rod of mass `$M$' and length `$2L$' is suspended at its middle by a wire. It exhibits torsional oscillations; If two masses each of `$m$' are attached at distance `$L/2$' from its centre on both sides, it reduces the oscillation frequency by 20\%. The value of ratio $m/M$ is close to:
\qn{2} Two concentric coplanar circular loops made of wire, with resistance per unit length $10\Omega$/m have diameters 0.2m and 2m. A time varying potential difference $(4+2.5t)$ volt is applied to the larger loop. The current in the smaller loop is $n\times10^{-10}$ A, determined the value of $n$.
\qn{3} Voltage rating of a parallel plate capacitor is 500 V. Its dielectric can withstand a maximum electric field of $10^6$ V/m. The plate area is $10^{-4}$m$^2$. What is the dielectric constant if the capacitance is 15pF ? (given $\epsilon_0=8.86\times10^{-12}$ C$^2$ m$^2$)
\qn{4} A thin sheet is inserted vertically inside a liquid of density $\rho$ and surface tension $T$. The liquid rises by height `$h$' close to sheet making angle of $30^\circ$ with sheet, as shown in the figure. The value of `$h$' in mm is \underline{\hspace{1cm}}. $[\rho=10^3$ kg/m$^3$, $T=0.81$ N/m, $g=10$ m/s$^2]$
\qfig{p2_phy_s2_q04_1}
\qn{5} A particle having mass $m$, charge $q$ is projected with velocity $v_0$ along y-axis in a region of uniform magnetic field $B_0$ which is outward and perpendicular to the plane of the paper as shown in the figure. The particle is continuously subjected to a frictional force which varies with velocity as $\vec F_r=-\alpha\vec v$, where $\alpha$ is a constant. Consequently, the particle moves on a spiral path till it comes to rest at point P. Find the x-co-ordinate (in metre) of point P. (Take $\alpha=10^{-3}$ kg/sec, $q=10^{-3}$C, $B_0=1$T, $v_0=1$ m/s, $m=20$ gm)
\qfig{p2_phy_s2_q05_1}
\qn{6} An inductor of 0.1 H and a box are connected to AC supply of 25 volts and $\omega=100$ rad/s. Over all power factor of circuit is 1 and the reading of AC ammeter ($I_{\text{rms}}$) is 5A. Find the power factor of box.
\qfig{p2_phy_s2_q06_1}
\sect{SECTION-III}
\qn{1} Three identical, parallel conducting plates A, B and C are placed as shown. Switches $S_1$ and $S_2$ are open and can connect A and C to earth when closed. A charge $Q=-3~\mu$C is given to B. If $S_1$ and $S_2$ both are closed together. Find the amount of charge in $\mu$C on inner surface of plate A.
\qfig{p2_phy_s3_q01_1}
\qn{2} A convex lens of focal length 1.5m is placed in a system of coordinate axis such that its optical centre is at origin and principal axis coinciding with the x-axis. An object and a plane mirror are arranged on the principal axis as shown in figure. Find the value of $d$ (in m) so that y-coordinate of final image (after refraction and reflection) is 0.3m. (Take $\tan\theta=0.3$)
\qfig{p2_phy_s3_q02_1}
\qn{3} Two parallel vertical metallic rails AB and CD are separated by 1m. They are connected at the two ends by resistances $R_1$ and $R_2$ as shown. A horizontal metallic bar PQ of mass 0.2 kg slides without friction, vertically down the rails under the action of gravity. There is uniform horizontal magnetic field of 0.6T perpendicular to plane of the rails. It is observed that when the terminal velocity attained, the power dissipated in $R_1$ and $R_2$ are 0.76 W and 1.2 W respectively. Find the terminal velocity of bar in m/s. ($g=9.8$ m/s$^2$)
\qfig{p2_phy_s3_q03_1}
\qn{4} The figure shows two identical parallel plate capacitors connected to a battery with switch S closed. The switch is now opened and the free space between the plates of the capacitors is filled with a dielectric of dielectric constant (or relative permittivity) 3. The ratio of the total electrostatic energy stored in both capacitors before and after the introduction of the dielectric is $\dfrac x5$. Then find the value of $x$
\qfig{p2_phy_s3_q04_1}
\qn{5} A particle is projected vertically upwards from the surface of the earth with a kinetic energy equal to $\dfrac13$ times the minimum kinetic energy needed to escape. If radius of the earth is 6400km. The maximum height attained by the particle from the surface of the earth is $\dfrac1N$ times the radius of earth. Determine the value on $N$.
\qn{6} A conducting rod of mass $m$ and length $\ell$ is connected by two identical springs as shown in the figure. Initially the system is in equilibrium. A uniform magnetic field of magnitude $B$ directed perpendicular to the plane of the paper outwards also exists in the region. If a current $I$ is switched on that passes from P to Q through the rod. Maximum elongation in the spring is $\dfrac{BI\ell}{NK}$. Determine the value of $N$.\par
[Given: $|mg|=|BI\ell|$]
\qfig{p2_phy_s3_q06_1}
\partsec{2}{CHEMISTRY}
\sect{SECTION-I}
\qn{1} Which one of the following compounds will give HVZ reaction?
\begin{qoptions}
\opt{A}{\qfigi{p2_chem_s1_q01_A}}
\opt{B}{\qfigi{p2_chem_s1_q01_B}}
\opt{C}{\qfigi{p2_chem_s1_q01_C}}
\opt{D}{\qfigi{p2_chem_s1_q01_D}}
\end{qoptions}
\qn{2} The reaction(s) leading to the formation of 1,3,5-trimethylbenzene is (are)
\begin{qoptions}
\opt{A}{\qfigi{p2_chem_s1_q02_A}}
\opt{B}{Me--C$\equiv$C--H $\xrightarrow[\text{873 K}]{\text{heated iron tube}}$}
\opt{C}{\qfigi{p2_chem_s1_q02_C}}
\opt{D}{\qfigi{p2_chem_s1_q02_D}}
\end{qoptions}
\qn{3} In which of the following I$^{\text{st}}$ is more stable than II$^{\text{nd}}$
\begin{qoptions}
\opt{A}{\qfigi{p2_chem_s1_q03_A}\ \qfigi{p2_chem_s1_q03_2}}
\opt{B}{\qfigi{p2_chem_s1_q03_B}}
\opt{C}{\qfigi{p2_chem_s1_q03_C}}
\opt{D}{None}
\end{qoptions}
\qn{4} Find the incorrect order in the following.
\begin{qoptions}
\opt{A}{He $<$ Ne $<$ Ar $<$ Kr $<$ Xe (boiling point)}
\opt{B}{He $>$ Ne $>$ Ar $>$ Kr $>$ Xe (ionization energy)}
\opt{C}{He $<$ Ne $<$ Ar $<$ Kr $<$ Xe (Solubility in water)}
\opt{D}{He $>$ Ne $>$ Ar $>$ Kr $>$ Xe (melting point)}
\end{qoptions}
\qn{5}
\qfig{p2_chem_s1_q05_1}
Identify the correct statement :-
\begin{qoptions}
\opt{A}{Salt A having $\text{S}_2\text{O}_3^{2-}$}
\opt{B}{White ppt. (D) is BaSO$_4$}
\opt{C}{B is H$_2$S}
\opt{D}{C form yellow ppt. with HgCl$_2$}
\end{qoptions}
\qn{6} Which reactions produce Chlorine?
\begin{qoptions}
\opt{A}{KMnO$_4$(aq.) $+$ HCl (conc.) $\to$}
\opt{B}{NaCl$_{\text{(aq.)}}$ (Brine) $\xrightarrow[\text{using Pt electrode}]{\text{Electrolysis}}$}
\opt{C}{NaOCl$_{\text{(aq.)}}$ $+$ HCl$_{\text{(aq.)}}$ $\to$}
\opt{D}{NaCl$_{\text{(s)}}$ $+$ H$_3$PO$_4$ (conc.) $\xrightarrow{\Delta}$}
\end{qoptions}
\sect{SECTION-II}
\qn{1} At what temperature (in K) is the following reaction at equilibrium, when all substances are at standard pressure? Assuming that entropies and enthalpies of reaction do not vary with temperature.\par
$\text{PCl}_{3(g)}+\text{Cl}_{2(g)}\rightleftharpoons\text{PCl}_{5(g)}$\par
\smallskip
{\footnotesize\setlength{\tabcolsep}{3pt}\renewcommand{\arraystretch}{1.4}\noindent\begin{tabular}{|c|c|c|}
\hline
 & $\Delta\text{H}^0_f$ (kJ/mole) & S$^\circ$ (J/mole-K)\\ \hline
PCl$_{3(g)}$ & $-290$ & 310\\ \hline
Cl$_{2(g)}$ & 0 & 220\\ \hline
PCl$_{5(g)}$ & $-375$ & 360\\ \hline
\end{tabular}\par}
\qn{2} 2.56gm of sulphur S$_8$(s) is taken which achieves equilibrium with its vapours according to reaction\par
S$_8$(s) $\rightleftharpoons$ 8S(g)\par
If vapours occupies 960 ml at 1atm, 300K, then degree of dissociation of S$_8$(s) will be -\par
[R = 0.08 atm L / mol-K] [Atomic mass : S = 32]
\qn{3} Number of compounds show Carbylamine test, or Isolynise test
\qfig{p2_chem_s2_q03_1}
\qfig{p2_chem_s2_q03_2}
\qfig{p2_chem_s2_q03_3}
\qn{4} Number of compound show more rate of EAS than benezene.
\qfig{p2_chem_s2_q04_1}
\qn{5} How many compound show Fehling and 2, 4 DNP test.\par
\noindent\qfigi{p2_chem_s2_q05_2}\hfill\qfigi{p2_chem_s2_q05_1}\hfill\qfigi{p2_chem_s2_q05_3}\par
\noindent\qfigi{p2_chem_s2_q05_5}\hfill\qfigi{p2_chem_s2_q05_4}\par
\noindent\qfigi{p2_chem_s2_q05_7}\hfill\qfigi{p2_chem_s2_q05_6}
\qn{6} The rate constant of the reaction $2\text{H}_2\text{O}_2\text{(aq)}\to2\text{H}_2\text{O(l)}+\text{O}_2\text{(g)}$ is $3\times10^{-3}$ min$^{-1}$. At what concentration of H$_2$O$_2$ (in M) the rate of the reaction will be $2\times10^{-4}$ M s$^{-1}$.
\sect{SECTION-III}
\qn{1} In a cyclic process ABCA consist of three processes. if $q_{AB}=-20$ J ; $q_{BC}=-40$ J ; $q_{CA}=30$ J \& $W_{AB}=80$ J ; $W_{BC}=-40$ J then find work done by system in process CA.\par
\textbf{\textit{Fill your answer as sum of digits (excluding decimal places) till you get the single digit answer.}}
\qn{2} Consider two reactions at $27^\circ$C\par
(I). A $\to$ B ; ( I order kinetics)\par
(II). X $\to$ Y; ( Zero order kinetics)\par
If both A \& X each having initial concentration 0.1M \& same initial half life period. If ratio of rate constants for I \& II reactions $\left(\dfrac{K_I}{K_{II}}\right)$ is `X' find value of $\left(\dfrac X2\right)$ ? [$\ln2=0.7$]
\qn{3} KF crystallizes in the NaCl type structure. If the radius of K$^+$ ion is $\dfrac{186}{\sqrt2}$ pm and F$^-$ ion is $\dfrac{214}{\sqrt2}$pm, what is the shortest distance between K$^+$ -- K$^+$ ions (in \AA\ in KF crystal) ?
\qn{4} Find the maximum number of identical S-O bond length in the anionic part of Na$_2$S$_2$O$_8$.
\qn{5} Find the number of moles of PH$_3$ formed when one mole of white phosphorus is treated with NaOH.
\qn{6} Find the moles of Na$_2$SO$_3$ required for the reduction of one mole of K$_2$Cr$_2$O$_7$ in acidic solution.
\partsec{3}{MATHEMATICS}
\sect{SECTION-I}
\qn{1} Which of the following function(s) is(are) surjective ?\par
[Note: $[m]$ and $\{m\}$ denotes greatest integer function and fractional part function respectively, and $D_\ell$ denotes the domain of the function $y=\ell(x)$]
\begin{qoptions}
\opt{A}{$f:D_f\to R$, $f(x)=\ell n\left(\tan(\pi[x])+\left|x^2+2x-3\right|\right)$}
\opt{B}{$g:D_g\to R$, $g(x)=\dfrac{x^2+2x-3}{x-1}$}
\opt{C}{$h:D_h\to R$, $h(x)=\ell n\left(\dfrac{1-x}{1+x}\right)$}
\opt{D}{$k:D_k\to R^+$, $k(x)=\sqrt{[x]+[-x]+1}+\sqrt{\{x\}+\{-x\}+1}$}
\end{qoptions}
\qn{2} There are $n$-students, namely $S_1,S_2,S_3,\ldots,S_n$ and $n$-chairs placed in a row. Let $Q_n$ denotes the number of ways of seating these $n$-students such that no two students having consecutive numbers sit together. Then
\begin{qoptions}
\opt{A}{$Q_3=0$}
\opt{B}{$Q_4=2$}
\opt{C}{$Q_5=14$}
\opt{D}{$Q_{21}>19Q_{20}$}
\end{qoptions}
\qn{3} Let $f$ be a continuous function on $R$ and satisfies $f(x)=e^x+\displaystyle\int_0^1e^xf(t)\,dt$, then which of the following is(are) correct?
\begin{qoptions}
\opt{A}{$f(0)<0$}
\opt{B}{$f(x)$ is decreasing function on $R$}
\opt{C}{$f(x)$ is an increasing function on $R$}
\opt{D}{$\displaystyle\int_0^1f(x)\,dx>0$}
\end{qoptions}
\qn{4} If $\displaystyle\int\dfrac{x+\sqrt[3]{x^2}+\sqrt[6]{x}}{x\left(1+\sqrt[3]{x}\right)}\,dx=\alpha\cdot x^\beta+\gamma\tan^{-1}(x^\delta)+C$, where $C$ is constant of integration, then which of the following is(are) correct? $\{2\alpha,3\beta,4\gamma,6\delta\in N\}$
\begin{qoptions}
\opt{A}{$\alpha=\dfrac23$}
\opt{B}{$\alpha\beta=1$}
\opt{C}{$(\alpha\beta+\gamma)=7$}
\opt{D}{Sum of all coefficients of even powers of $x$ in the expansion of $\left(\alpha\beta+\dfrac x\delta\right)^\gamma$ is equal to $\dfrac{(\gamma+1)^\gamma+(\gamma-1)^\gamma}{2}$.}
\end{qoptions}
\qn{5} The value of the definite integral $\displaystyle\int_{-1}^0\left(\sqrt[3]{x^4-2x^2}-\sqrt{1-\sqrt{x^3+1}}\right)dx$ lies in the interval
\begin{qoptions}
\opt{A}{$\left(\dfrac{-\pi}2,\dfrac\pi2\right)$}
\opt{B}{$(-\pi,\pi)$}
\opt{C}{$\left(\dfrac\pi4,\dfrac\pi2\right)$}
\opt{D}{$\left(\dfrac{-\pi}3,\dfrac\pi3\right)$}
\end{qoptions}
\qn{6} If\par\noindent\hfill$F(x)=\displaystyle\int e^x\left(\int_0^xx\sin x^2dx+\int_0^x\left(\sin x^2+2x^2\cos x^2\right)dx\right)dx$\hfill\null\par\noindent$F(0)=\dfrac{-1}2$ then which of the following statement(s) must be true.
\begin{qoptions}
\opt{A}{$F\left(\sqrt{\dfrac{4\pi}7}\right)<0$}
\opt{B}{$F'(0)=\dfrac{-1}2$}
\opt{C}{$F''(0)+F'''(0)=-1$}
\opt{D}{$F(0)-F'(0)=\dfrac12$}
\end{qoptions}
\sect{SECTION-II}
\qn{1} $f(x)$ is continuous and bijective function on $R$. If $\forall\ a\in R$, the area bounded by $y=f(x)$, $x=3-a$, $x=3$ and x-axis is equal to area bounded by $y=f(x)$, $x=3+a$, $x=3$ and x-axis. Then the value of $k$ if $\displaystyle\int_{-\lambda}^{\lambda}f^{-1}(x)\,dx=k\lambda$ is
\qn{2} If $f'(x)=\phi(x)$ and $\phi'(x)=f(x)\ \forall\ x\in R$, also $f(3)=5$ and $f'(3)=4$, then value of $\left(f^2(10)-\phi^2(10)\right)$ is
\qn{3} Consider a twice differentiable function $y=f(x)$ satisfying $f(x)+f''(x)=2f'(x)$ where $f(0)=0$, $f(1)=e$ then find $\displaystyle\int_0^1f(x)\,dx$
\qn{4} Given $p,q\in\{0,1,2,3,4,\ldots\ldots\ldots,19,20\}$. Consider the system of equations\par
\hspace*{1.5em}$x+y+z=4$\par
\hspace*{1.5em}$2x+y+3z=6$\par
\hspace*{1.5em}$x+2y+pz=q$\par
Let\ \ L: denotes number of ordered pairs $(p,q)$ so that the system of equations has a unique solution,\par
\hspace*{1.2em}M: denotes number of ordered pairs $(p,q)$ so that the system of equations has no solution and\par
\hspace*{1.2em}N: denotes number of ordered pairs $(p,q)$ so that the system of equations has infinite solutions.\par
Find $(L+M-N)$.
\qn{5} Let $f$ be a differentiable function satisfying $f'(x)=f^2(x)$ and $f(0)=-1$.\par
If $\dfrac{\sqrt2}{\pi}\displaystyle\int_0^1\left(\dfrac{(1-x)}{f(x)(x^4+1)}+\dfrac{2}{x^4+1}\right)dx=\dfrac pq$ where $p,q\in N$ then find least value of $(p+q)$.
\qn{6} If area of the region whose upper left boundary is enclosed by the curve $f(x)=\sin x+x+2$, lower left by the line $x+y=2$, lower right by $f^{-1}(x)$ and upper right by the line $x+y=2\pi+2$ is $k(\pi+1)$ then find the value of $k$
\sect{SECTION-III}
\qn{1} Let $y=f(x)$ be the curve passing through $M(0,2)$ such that the sum of coordinates of any point $P(x,y)$ on the curve exceeds the slope of tangent to the curve at that point by 5. Find the number of distinct real solutions of the equation $f'(x)=0$.
\qn{2} Let $\alpha(t)$ and $\beta(t)$ be differentiable functions on $R$ such that $\alpha(0)=2$ and $\beta(0)=1$. If $\alpha(t)+\beta'(t)=1$ and $\alpha'(t)+\beta(t)=1$ for all $t\in[0,\infty)$, then the value of $\alpha(\ell n\,2)$ is expressed in the lowest form as $\dfrac pq$. Find the value of $(p-q)$.
\qn{3} A square matrix of size $n\times n$ is called an orthogonal matrix iff $AA^T=A^TA=I_n$, $n\in N$.\par
If $A=\begin{bmatrix}a&b&c\\c&a&b\\b&c&a\end{bmatrix}$ is an orthogonal matrix, then magnitude of $(a+b+c)$ is equal to
\qn{4} Let $A$ and $B$ be orthogonal matrices of order $3\times3$ such that $X=AB^T$ and $Y=BA^T$, then the value of $\text{tr}(XY)$ is equal to -
\qn{5} Let matrix $A=\begin{bmatrix}\alpha^2&6&8\\3&\beta^2&9\\4&5&\gamma^2\end{bmatrix}$ and $B=\begin{bmatrix}2\alpha&3&5\\2&2\beta&6\\1&4&2\gamma-3\end{bmatrix}$. If $\text{tr}(A)=\text{tr}(B)$ then value of $\left(\dfrac1\alpha+\dfrac1\beta+\dfrac1\gamma\right)$ is equal to (where $\text{tr}(A)$ denotes trace of $A$ and $\alpha,\beta,\gamma$ are non-zero real numbers)
\qn{6} $W_1$, $W_2$ and $W_3$ are three ticket windows in a cinema hall. 12 persons ($A_1,A_2,\ldots\ldots,A_{12}$) are standing for the tickets in a queue of 4, 3 and 5 number of persons infront of $W_1$, $W_2$ and $W_3$ respectively. $A_1$ and $A_2$ want to get their tickets from $W_1$ and $A_3$, $A_4$ and $A_5$ want to get their tickets from $W_3$. If the number of ways in which all are getting their tickets such that $A_1$, $A_2$, $A_3$, $A_4$ and $A_5$ are getting their tickets as early as possible is $k(7!)$, then find the value of $\left[\dfrac k3\right]$.\par
(Assuming each person is getting exactly one ticket at a time).
[Note: $[y]$ denotes greatest integer less than or equal to $y$.]
\end{multicols}
\end{document}
